% =====================================================================
%  example.tex
%  Example paper built on the eupvsec.cls class.
%
%  All layout rules (fonts, margins, columns, heading/caption style)
%  live in eupvsec.cls - do not change them here. Only edit the
%  content below, and cite sources from refs.bib with \cite{key}.
% =====================================================================
\documentclass{eupvsec}

\begin{document}

\title{INSTRUCTIONS FOR THE PREPARATION OF PAPERS\\FOR THE EU PVSEC 2026}
\paperauthors{Author(s)}
\paperaffiliation{Company / Institute(s)}
\paperaddress{Address(es)}

\paperabstract{%
These notes include important information on how to
prepare and submit your paper. Read the notes carefully and follow
them precisely. \textbf{Your paper must be written in English}. The
layout of your paper should be exactly the same as this master
document. \textbf{In order to prepare your layout, save this
document with a new name (use your session code as the name for the
document, e.g.\ 4AV.1.20) and use it as a template. Substitute the
following text with the text of your paper without changing the
layout} (see Section~3). Do not change font type and size, line
spacing, page margins and the structure of this document.
\textbf{Do not insert page numbers and page headers/footers}.
Ignoring these instructions might cause problems and delays in the
publication of the Conference Proceedings. \textbf{In case your
abstract was selected for participation in the peer review process
for publication in a Scientific Journal} (PiP, Solar RRL, EPJ
Photovoltaics or AESR), as well as your submission to the journal,
\textbf{you can also publish your Short Paper in the conference
proceedings (N.B. it is essential that the contents of both papers
differ by at least 50\%)}. Published papers will include a note to
indicate that the Full Paper was submitted to the selected Journal.
The Conference Proceedings will be online searchable, online citable
and available for full permanent free access and download on the
EU PVSEC Proceedings website (www.eupvsec-proceedings.com). If you
have any questions, please do not hesitate to contact us
(see Section~6.3).}

\paperkeywords{Keywords: three to five keywords in order of importance}

\maketitle

% =====================================================================
%   BODY  (two columns from here on)
% =====================================================================

\section{Deadlines and delivery}

Your original paper has to be submitted online between \textbf{07th --
17th September 2026 through the corresponding author\'s user area
(\href{https://userarea.eupvsec.org}{userarea.eupvsec.org})}. Kindly
note that \underline{only Corresponding Authors of each submission}
may upload final manuscripts. It means that the corresponding author
is the only author from each paper that is able to complete the
submission (as is the case for abstract submission). If the manuscript
is not available during this period, your paper \textbf{cannot be
published} in the Conference Proceedings.

The document must be submitted in \textbf{Adobe Acrobat PDF} format.

To upload the paper, follow the step-by-step procedures provided in
the user area.

In addition to the manuscript, \textbf{complete the mandatory
electronic Copyright Transfer Agreement as one of the steps of the
online submission.}

\textbf{Please make sure that the paper you submit is the final
version.}

Your paper can only be published if it was presented during the
conference and at least one of the authors (presenter) is duly
registered. The \textbf{Conference Registration Area} is available on
the website
\href{https://www.eupvsec.org/index.php/registration}{Registration (eupvsec.org)}.

\section{Preparing the manuscript}

\subsection{Volume of paper}

The complete paper should not exceed \textbf{1~MB including
illustrations}. This is a very good proven capacity for final papers.
Note that there is no limit in the number of pages. An adequate number
is 3 to 6 pages.

\subsection{Organisation of the paper}

The \textbf{title} of the paper should be informative and concise. It
should be followed by the author(s) name(s) -- listing the principal
author first --, organisation, address, telephone and email address.

Precede the body of your paper with an \textbf{abstract} of 100 to 200
words, giving a brief account of the most relevant aspects of your
paper. Please avoid using symbols, graphics and formulas in this part
of the document as this section will appear as plain text on the
EU PVSEC 2026 Proceedings website.

Next, in order of importance, include three to five most relevant
\textbf{keywords} in your paper, separating them with commas.

The body of the text must be in \textbf{two columns}. Number each
heading using decimal numbering. Follow the layout specifications in
Section~3 below.

\section{Layout specifications}

\textbf{The layout of your paper should have exactly the same format
as these notes.}

Select the printer option ``Acrobat Distiller'' or Adobe PDF
(version~8.0 and up) before starting to work on this document in order
to avoid accidental displacement of layout elements afterwards when
converting the Word file into PDF format. Please abstain from freeware
products which are not 100\% compatible with Adobe PDF. Should you not
have the above technical requirements, please contact us (see 6.3).

\subsection{Fonts and type size}

Font type: Times New Roman. Font size: 9pt. Line spacing: single. Text
alignment: justified left and right. Captions should be the same font
and size as the typeface used for the text. Make sure that
illustrations are clear and easy to read. Please do not use any other
font than Times New Roman.

\subsection{Page size}

Page size must be A4 (210mm x 297mm).

Margins: top: 32mm; bottom: 19mm; left: 25mm; right: 25mm. Text or
images can't be outside these margins. The body of the text must be
in two equal columns. The width of the columns should be 73.6mm.

\subsection{Layout of the text}

Begin at the top of the first page with the title of the paper, in
capital letters, bold and centered. Leave two blank lines between the
title and the name of the author(s), and list the surname preceded by
the first name. On the following lines, give the name of the company
or institute, wherever applicable, and the full address. When several
authors prepare a paper, the name of the principal author should
appear first and the name of every institute should be easy to depict.
Next, leave two blank lines and then type an abstract of no more than
200 words (indent the whole block as at the top of this document). At
the end of the abstract, give your list of keywords. Leave another two
blank lines between the abstract and the body of the paper, which
should be typed in two columns.

\subsection{Headings}

Leave two blank lines before each section and one blank line before
the heading of each sub-section. Headings and sub-headings should be
numbered (e.g.\ 3, 3.1, 3.2). There should be no blank line after the
title of the sub-sections, only an indentation to indicate the
beginning of a paragraph. Section headings should be in capital
letters, sub-section headings in upper and lower case. Headings should
be normal text -- not underlined or in bold.

\section{Additional components}

\subsection{Illustrations}

Illustrations (drawings, graphs, photographs, etc.) should not exceed
50\% of the whole paper and should be placed near their citation.
Accepted formats are: \textbf{.jpg, .png} and \textbf{.gif}. Please
note that vector graphics need to be converted into bit-mapped
graphics.

Illustrations should have a resolution of \textbf{300~dpi}, using
simple colours (Standard + RGB) and be placed at \textbf{100\% scale}
(i.e.\ if an illustration covers the full column width, it should be
of approx.\ 860 pixels). \textbf{Illustrations should not be taken
from previously-printed material.}

Illustrations should have the layout in line with the text.

Figures must be numbered (e.g.\ Fig.~1) in the text and in the figure
caption. Captions should be clear enough to allow comprehension of the
illustration without reference to the text.

Graphs must not be imported from Excel, but be inserted as a picture
(.jpg, .png or .gif). Please use simple contrasting colours (Standard
+ RGB) instead of fill patterns. See Fig.~1 for good/bad example.

Please make sure that the \textbf{PDF} file size of the final document
does not exceed \textbf{1~MB}. If your document size should be larger
than this, verify the resolution and scale down the inserted images.

% --- Example figure -------------------------------------------------
\begin{figure}[t]
  \centering
  \fbox{\parbox[c][2.6cm][c]{\columnwidth}{\centering
    \small Good example with contrasting colours\\(placeholder graphic)}}
  \caption{Clear line drawings are essential.}
  \label{fig:example}
\end{figure}
% ---------------------------------------------------------------------

\subsection{Tables}

Tabular presentation of data is an easy way of condensing many items.
Refer to tables by using Roman numerals. All tables must be numbered
and captioned (see Table~\ref{tab:swr}).

\begin{table}[t]
  \centering
  \caption{Standing waves ratio}
  \label{tab:swr}
  \begin{tabular}{lccc}
    \hline
     & 110\,MHz & 120\,MHz & 130\,MHz \\
    \hline
    Standing waves ratio   & 1.86 & 1.73 & 1.40 \\
    Reflection coefficient & 0.32 & 0.28 & 0.16 \\
    Phase                  & 3.78 & 3.89 & 4.02 \\
    \hline
  \end{tabular}
\end{table}

\subsection{References}

References should not appear as footnotes but should be gathered
together at the end of the text. When referring to them in the text,
type the corresponding reference number in brackets, e.g.\
\cite{smith2002}. References should be numbered as shown at the end of
this document -- here they are pulled automatically from
\texttt{refs.bib} with \verb|\cite{}| and the \texttt{cite} package
(see \cite{smith2002,boer1996}).

\section{Special requests for the PDF version}

\subsection{Printer}

Select Adobe PDF (version~8.0 and up) as the printer option to convert
your Word file into PDF version (1.6 or 1.7).

Go to printer properties and click on ``Acrobat PDF settings'', then
choose as conversion settings: ``Press'' quality.

\subsection{Password / Security}

\textbf{Please do not use passwords nor security when making your PDF
file.} You may rest assured that we do not make any changes in your
manuscript. As the header and page numbering will be added by us, we
need to ensure that your PDF file is not password protected.

\section{Other points}

\subsection{Permissions}

Authors are fully responsible for their papers. They must take the
necessary steps to obtain permission for using any material that
might be protected by copyright.

\subsection{Copyright}

Please be aware that on submission of your manuscript, you transfer
the \textbf{copyright} to WIP, publisher of the proceedings.

\subsection{Further information}

\textbf{Papers must be submitted online between 07th to 17th September
to the user area} of corresponding authors of each paper at
\url{https://userarea.eupvsec.org/}. All requests regarding the
preparation of the paper should be addressed to:

\medskip
\noindent
WIP Renewable Energies\\
attn. EU PVSEC Programme Secretariat\\
Tel. +49-89-720 12 -735, Fax +49-89-720 12 791\\
E-mail: \href{mailto:pv.manuscripts@wip-munich.de}{pv.manuscripts@wip-munich.de}
\medskip

\section{Summary of these notes}

\begin{itemize}
  \item Use this layout as a guide to write your paper. Write your
        paper directly on this document and save it with the name of
        the session code.
  \item Select Adobe PDF (version~8.0 and up) before starting to work
        on this document.
  \item Write your paper in English, A4 format, Font Times New Roman,
        9~pt.
  \item Use contrasted pictures/illustrations placed at 100\% (do not
        resize them in Microsoft Word), 300~dpi with Standard + RGB
        colors in .jpg, .png or .gif format.
  \item The size of the final document in PDF cannot be larger than
        1~MB.
  \item Submit your paper online between \textbf{07th to 17th
        September to the user area} of corresponding authors of each
        paper at \url{https://userarea.eupvsec.org}, in Adobe Acrobat
        PDF format.
  \item Complete the mandatory electronic Copyright Transfer Agreement
        as one of the steps of the online submission.
\end{itemize}

\section{Author's checklist for delivery of the final manuscript}

Before submitting your manuscript at the Conference please make sure
that:

\begin{itemize}
  \item Your paper follows the layout specified in these guidelines
  \item Your paper does NOT contain header or page numbering
  \item The digital version of your manuscript is uploaded in the
        user area in Adobe Acrobat PDF format
  \item The digital files are named with the name of the session
        reference code
  \item Your PDF is NOT password protected
\end{itemize}

If you need any further assistance, please contact us.

Thank you very much for your kind attention.

The Publisher

% =====================================================================
%   REFERENCES  (generated by BibTeX from refs.bib)
% =====================================================================
\bibliographystyle{unsrt}
\bibliography{refs}

\end{document}
